\documentclass[12pt,ngerman,parskip=half]{scrreprt} % draft für keine Bilder
\usepackage{blindtext}

\usepackage[T1]{fontenc}
\usepackage{xcolor}
\usepackage{graphicx}
\usepackage{booktabs}
\usepackage{paralist}
\usepackage{listings}
\usepackage{csquotes}
\usepackage{lmodern}
\usepackage{here}
\usepackage{amsmath}

\makeatletter
\newcommand*{\avg}{\mathop{\operator@font avg}}
\makeatother

\usepackage{babel}
\usepackage{microtype}

\author{Uwe Ziegenhagen}
\title{Mein erstes \LaTeX-Dokument}

\newcommand{\person}[1]{\textcolor{red}{\textsc{#1}}}

\newcommand{\pperson}[2]{\textbf{#1}~\textsc{#2}}

\usepackage{hyperref}
\hypersetup{
    bookmarks=true,                     % show bookmarks bar
    unicode=false,                      % non - Latin characters in Acrobat’s bookmarks
    pdftoolbar=true,                        % show Acrobat’s toolbar
    pdfmenubar=true,                        % show Acrobat’s menu
    pdffitwindow=false,                 % window fit to page when opened
    pdfstartview={FitH},                    % fits the width of the page to the window
    pdftitle={My title},                        % title
    pdfauthor={Author},                 % author
    pdfsubject={Subject},                   % subject of the document
    pdfcreator={Creator},                   % creator of the document
    pdfproducer={Producer},             % producer of the document
    pdfkeywords={keyword1, key2, key3},   % list of keywords
    pdfnewwindow=true,                  % links in new window
    colorlinks=true,                        % false: boxed links; true: colored links
    linkcolor=blue,                          % color of internal links
    filecolor=blue,                     % color of file links
    citecolor=blue,                     % color of file links
    urlcolor=blue                        % color of external links
}

\begin{document}
\maketitle

\tableofcontents

\listoffigures

\listoftables

\chapter{Themenüberblick}

\section{Einführung}
\subsection{Literaturüberblick}
\subsubsection{Literatur vor 1900}

Hallo, ich {\tiny bin} ein {\Huge großer} Text. \person{Albert Einstein} \pperson{Albert}{Einstein} 

Siehe dazu auch Abschnitt \ref{sec:hauptteil} auf Seite \pageref{sec:hauptteil}.

\person{Albert Einstein}\footnote{Bekannter Physiker}\marginpar{Albert Einstein war ein bekannter Physiker, der sich mit der Relativitätstheorie beschäftigte.}

Albert Einstein sagte einmal: \enquote{Glaubt nicht alles, \enquote{Glaubt nicht alles, was ihm im Internet lest!} was ihm im Internet lest!}

Albert Einstein sagte einmal: \glqq Glaubt nicht alles, was ihm im Internet lest!\grqq

Siehe Abbildung \ref{fig:katze}  auf Seite \pageref{fig:katze}

Siehe Abbildung \ref{fig:katze2}  auf Seite \pageref{fig:katze2}

\textmd{Uwe Ziegenhagen}

\textbf{Uwe Ziegenhagen}

\textcolor{red}{Ich bin roter Text}

Hallo, ich bin \textit{Text} \textit{\textbf{Text}} \textbf{\textit{Text}}   ein \texttt{Typewriter} \textbf{fettgedruckter} Text

\blindtext \clearpage

\begin{itemize}
	\item Hallo
	\item ich bin
	\item eine 
	\item itemize
	\item Umgebung
	\item in LaTeX
\end{itemize}

\begin{compactitem}[\textcolor{red}{$\Rightarrow$}]
	\item Hallo
	\item ich bin
	\item eine 
	\item itemize
	\item Umgebung
	\item in LaTeX
\end{compactitem}

\begin{enumerate}
	\item Hallo
	\item ich bin
	\item eine 
	\item itemize 
	\item Umgebung
	\item in LaTeX
\end{enumerate}

\begin{compactenum}[1.] \setcounter{enumi}{2} 
	\item Hallo
	\item ich bin
	\item eine 
	\item itemize 
	\item Umgebung
	\item in LaTeX
\end{compactenum}

\begin{description}
\item[Apfel] ein Obst
\item[Birne] auch ein Obst
\item[Tomate] eher Gemüse
\end{description}

\begin{compactdesc}
\item[Apfel] ein Obst
\item[Birne] auch ein Obst
\item[Tomate] eher Gemüse
\end{compactdesc}


\blindtext

\blindtext

\subsubsection{Literatur nach 1900}

\blindtext

\blindtext

\blindtext

\section{Hauptteil}\label{sec:hauptteil}

\blindtext[1]


\section{Schlussworte}

\blindtext[10]

\begin{tabular}{|l|r|c|p{70mm}|} \hline
\textbf{1. Spalte} & \textbf{2. Spalte} & \textbf{3. Spalte} & \textbf{4. Spalte} \\ \hline \hline
Hallo Welt & 1.2345 & 123 & Hallo, ich bin ein kleiner Text. \\ \hline
Hallo Welt & 1.2345 & 123 & Hallo, ich bin ein kleiner Text. \\ \hline
Welt & 1.23 & 465465465 & Hallo \\ \hline
Hallo Welt & 1.2345 & 123 & Hallo, ich bin ein kleiner Text. \\ \hline
Hallo Welt & 1.2345 & 123 & Hallo, ich bin ein kleiner Text. \\ \hline
\end{tabular}\vspace*{2cm}

\begin{center}
\begin{tabular}{lrrp{60mm}} \toprule[2pt]
\textbf{1. Spalte} & \textbf{2. Spalte} & \textbf{3. Spalte} & \textbf{4. Spalte} \\ \cmidrule[1pt](rl){1-4}
Hallo Welt & 1.2345 & 123 & Hallo, ich bin ein kleiner Text. \\ 
Hallo Welt & 1.2345 & 123 & Hallo, ich bin ein kleiner Text. \\ 
Welt & 1.2300 & 465465465 & Hallo \\ 
Hallo Welt & 1.2345 & 123 & Hallo, ich bin ein kleiner Text. \\ 
Hallo Welt & 1.2345 & 123 & Hallo, ich bin ein kleiner Text. \\ \bottomrule
\end{tabular}
\end{center}

\begin{center}
\Large Hallo DLR!
\end{center}

\chapter{Bilder}

\includegraphics[width=5cm]{Katze2.jpg}

\includegraphics[width=\textwidth]{Katze2.jpg}

\begin{center}
\includegraphics[width=0.75\textwidth]{Katze2.jpg}
\end{center}

\includegraphics[height=\textheight,width=\textwidth]{Katze2.jpg}

ABC\hspace{10mm}DEF

ABC\hspace{1cm}DEF

ABC\hspace{10pt}DEF

ABC\hspace{10pc}DEF

ABC\hspace{10ex}DEF

ABC\hspace{10em}DEF



\begin{center}
\includegraphics[width=0.75\textwidth]{./Bilder/Katze/Katze2.jpg}
\end{center}

\chapter{Gleitobjekte}

\blindtext

\blindtext

% t = top, Anfang der Seite
% b = bottom, Ende der Seite
% h = here, da, wo das Bild eingefügt wurde
% ! = Sei entspannt, LaTeX, wo Du das Bild hinpackst
% H => über \usepackage{here} strenger als das kleine h

\begin{figure}[h] % tbhH!
\begin{center}
\includegraphics[width=0.75\textwidth]{./Bilder/Katze/Katze2.jpg}
\caption{Unsere Miezekatze}\label{fig:katze}
\end{center}
\end{figure}

Siehe Abbildung \ref{fig:katze}  auf Seite \pageref{fig:katze}

\blindtext

\blindtext

\blindtext

\begin{table}
\begin{center}
\caption{Meine schöne Tabelle}\label{tab:nice}
\begin{tabular}{lrrp{60mm}} \toprule[2pt]
\textbf{1. Spalte} & \textbf{2. Spalte} & \textbf{3. Spalte} & \textbf{4. Spalte} \\ \cmidrule[1pt](rl){1-4}
Hallo Welt & 1.2345 & 123 & Hallo, ich bin ein kleiner Text. \\ 
Hallo Welt & 1.2345 & 123 & Hallo, ich bin ein kleiner Text. \\ 
Welt & 1.2300 & 465465465 & Hallo \\ 
Hallo Welt & 1.2345 & 123 & Hallo, ich bin ein kleiner Text. \\ 
Hallo Welt & 1.2345 & 123 & Hallo, ich bin ein kleiner Text. \\ \bottomrule
\end{tabular}
\end{center}
\end{table}

\blindtext

\blindtext


\begin{center}
\includegraphics[width=0.75\textwidth]{./Bilder/Katze/Katze2.jpg}
\captionof{figure}{Unsere Miezekatze}\label{fig:katze2}
\end{center}


\chapter{Mathematik}

Ich bin eine Formel $a^2+b^2=c^2$ im laufenden Text in TeX-Notation.

$$a^2+b^2=c^2$$

Ich bin eine Formel \(a^2+b^2=c^2\) im laufenden Text in LaTeX-Notation.

\[a^2+b^2=c^2\]

\begin{equation}
- \frac{p}{2} \pm \sqrt{ \left(\frac{p}{2}\right)^2 - q }
\end{equation}

\begin{equation}
c_{a_b}^{2^2} \cdot \cdots \dots \ldots \int_{i=0}^{\infty} \not= \sum_{x}^{y}
\end{equation}


\begin{equation}
\overbrace{a^2 + b^2}^{\text{Hallo Welt}}_{456} =\underbrace{c^2 + d^2}_{\mbox{Hallo Welt}}
\end{equation}

\begin{equation}
sin \not= \sin \cos \tan \cosh \avg
\end{equation}

Siehe Formel \ref{eq:test}

\begin{eqnarray}
y  &=& d \\
  &=& d + e\\\label{eq:test}
  &=& d + \cos x \\
\end{eqnarray}



\end{document}

